\documentclass[../DoAn.tex]{subfiles}
\begin{document}

Chương này trình bày bức tranh toàn cảnh về sự gia tăng của các mối đe dọa
bảo mật ứng dụng web trong bối cảnh chuyển đổi số hiện nay, từ đó làm rõ tính
cấp thiết của bài toán phòng thủ ở tầng ứng dụng. Nội dung chương xác định
các hạn chế của các giải pháp tường lửa ứng dụng web hiện tại, đề xuất mục tiêu
cụ thể của đề tài, và giới thiệu định hướng giải pháp kỹ thuật mà đồ án theo đuổi.

% ====================================================================
\section{Đặt vấn đề}
\label{section:1.1}

Trong thập kỷ vừa qua, ứng dụng web đã trở thành hạ tầng trung tâm của nền
kinh tế số. Theo báo cáo của Verizon Data Breach Investigations Report (DBIR)
năm 2023, hơn 74\% các sự cố vi phạm dữ liệu có liên quan đến yếu tố con
người hoặc lỗ hổng ở tầng ứng dụng, trong đó tấn công vào ứng dụng web chiếm
vị trí đứng đầu trong các vectơ xâm nhập được ghi nhận. Tổ chức OWASP (Open
Web Application Security Project) hàng năm công bố danh sách OWASP Top 10,
tập hợp mười nhóm lỗ hổng nghiêm trọng và phổ biến nhất tại tầng ứng dụng.
Trong phiên bản cập nhật gần nhất (2021), các nhóm như SQL Injection (SQLi),
Cross-Site Scripting (XSS), Server-Side Request Forgery (SSRF), và các dạng
tấn công tiêm nhiễm lệnh (Command Injection) vẫn tiếp tục duy trì thứ hạng cao,
phản ánh thực trạng rằng chúng không chỉ phổ biến mà còn liên tục được khai thác
thực tế.

Trước áp lực đó, Tường lửa Ứng dụng Web (Web Application Firewall -- WAF)
ra đời và đóng vai trò như một tầng phòng thủ chuyên biệt, đặt giữa người dùng
và ứng dụng backend. Khác với tường lửa mạng thông thường vận hành ở tầng
3 và 4 của mô hình OSI, WAF kiểm tra nội dung thực sự của từng yêu cầu HTTP
ở tầng 7, cho phép phát hiện các hành vi tấn công ẩn bên trong payload, tham số
URL, tiêu đề HTTP hay nội dung body của yêu cầu. Đây là tầng kiểm soát mà
tường lửa mạng truyền thống không thể tiếp cận.

Tuy nhiên, giải pháp WAF dựa thuần túy trên luật (rule-based) -- hiện là phương
pháp thống trị trong hầu hết các triển khai thực tế -- bộc lộ những hạn chế cơ bản
khó khắc phục. Thứ nhất, bộ luật được xây dựng dựa trên các mẫu (pattern) tấn
công đã biết, chủ yếu thông qua biểu thức chính quy (regular expression). Kẻ tấn
công dày dạn kinh nghiệm có thể vượt qua các pattern này bằng kỹ thuật làm mờ
(obfuscation): mã hóa nhiều lớp (double URL encoding), thay đổi hoa thường kết
hợp với chèn ký tự trắng, dùng biểu thức thay thế tương đương về ngữ nghĩa nhưng
khác về cú pháp. Chẳng hạn, câu truy vấn \texttt{UNION SELECT} có thể được biến
thể thành \texttt{uNiOn\%09SELECT} hoặc \texttt{UN\%49ON SELECT} để qua mặt
các biểu thức chính quy kiểm tra case-sensitive. Thứ hai, hành vi của các rule
trong thực tế mang tính nhị phân (block hoặc allow), thiếu cơ chế đánh giá xác
suất theo ngữ cảnh toàn bộ yêu cầu. Một yêu cầu có thể khớp một phần với nhiều
rule nhưng không khớp hoàn toàn với bất kỳ rule đơn lẻ nào -- trường hợp này
rất dễ bị bỏ qua nếu không có cơ chế tích lũy điểm (anomaly scoring). Thứ ba,
bộ luật phải được cập nhật thủ công khi xuất hiện biến thể tấn công mới, tạo ra
khoảng thời gian chờ (detection gap) mà hệ thống hoàn toàn mù quáng trước các
kỹ thuật tấn công zero-day hoặc mới được công bố.

Ở chiều ngược lại, các giải pháp WAF thương mại trên nền tảng đám mây như
Cloudflare WAF hay AWS WAF giải quyết được một phần vấn đề cập nhật luật nhờ
đội ngũ nghiên cứu bảo mật chuyên trách và khả năng học từ lưu lượng toàn cầu.
Tuy nhiên, chúng lại phát sinh các hạn chế khác: toàn bộ lưu lượng ứng dụng,
bao gồm dữ liệu nhạy cảm của người dùng và khách hàng, phải đi qua hạ tầng
của bên thứ ba, đặt ra những lo ngại nghiêm trọng về quyền riêng tư và tuân thủ
quy định (GDPR, Nghị định 13/2023/NĐ-CP tại Việt Nam). Chi phí vận hành tỷ
lệ thuận với lưu lượng làm cho các giải pháp này khó tiếp cận với các tổ chức vừa
và nhỏ. Quan trọng hơn, tính chất ``hộp đen'' của các dịch vụ này khiến đội ngũ
bảo mật nội bộ mất đi khả năng kiểm soát, kiểm tra, và tinh chỉnh hành vi của
WAF theo đặc thù riêng của ứng dụng.

Trong bối cảnh đó, nhu cầu về một giải pháp WAF triển khai tại chỗ (on-premise),
mã nguồn mở, hiệu năng cao, có khả năng xử lý các yêu cầu ở tầng ứng dụng với
độ trễ tối thiểu đồng thời tích hợp thêm một lớp phân tích thông minh để xử lý
các trường hợp mơ hồ mà bộ luật cứng không thể giải quyết dứt khoát -- là nhu
cầu chính đáng và chưa có giải pháp mã nguồn mở nào đáp ứng thỏa đáng. Hệ
thống như vậy, nếu được xây dựng và mở rộng đúng hướng, sẽ đem lại lợi ích
cho nhiều đối tượng: từ các doanh nghiệp cần kiểm soát toàn bộ hạ tầng bảo mật
của mình, đến các đơn vị nghiên cứu muốn thử nghiệm và đánh giá các kỹ thuật
phòng thủ mới trong môi trường có thể tái lập (reproducible environment).

% ====================================================================
\section{Mục tiêu và phạm vi đề tài}
\label{section:1.2}

Để xác định mục tiêu và phạm vi phù hợp, đề tài trước hết tiến hành khảo sát
và đánh giá tổng quan các giải pháp WAF mã nguồn mở và thương mại đang được
sử dụng phổ biến nhất hiện nay.

ModSecurity kết hợp bộ luật OWASP Core Rule Set (CRS) là giải pháp WAF
mã nguồn mở lâu đời nhất và được triển khai rộng rãi nhất trong thực tế. ModSecurity
hoạt động như một module nhúng vào máy chủ web (Apache, Nginx, IIS) và xử lý
lưu lượng thông qua năm pha kiểm tra (connection, request headers, request body,
response headers, response body). Bộ luật OWASP CRS cung cấp hơn 200 rule
bao phủ hầu hết các nhóm tấn công trong OWASP Top 10. Dù vậy, ModSecurity
gặp phải những hạn chế đáng kể trong vận hành thực tế: ngưỡng mặc định của CRS
tạo ra tỷ lệ false positive cao trên các ứng dụng có payload phức tạp (API JSON,
tham số dài), đòi hỏi quá trình tinh chỉnh tốn công sức trước khi đưa vào sản xuất.
Kiến trúc module nhúng vào web server làm giới hạn khả năng triển khai độc lập
và tái cấu hình động mà không cần khởi động lại dịch vụ. Ngoài ra, bản thân
ModSecurity không có giao diện quản trị tích hợp, buộc đội ngũ vận hành phải
kết hợp với các công cụ bên thứ ba như Kibana hay Grafana để có khả năng giám
sát đầy đủ.

Cloudflare WAF và AWS WAF đại diện cho hướng tiếp cận đám mây, cung cấp
giao diện quản trị trực quan, khả năng cập nhật luật tức thời và tích hợp sẵn mạng
phân phối nội dung (CDN) để giảm thiểu các cuộc tấn công DDoS volumetric.
Cloudflare đặc biệt nổi trội với khả năng học từ lưu lượng toàn cầu và các quy
tắc dựa trên máy học (ML-based rules) bổ sung cho tập luật truyền thống. Tuy
nhiên, như đã phân tích ở phần \ref{section:1.1}, cả hai giải pháp đều yêu cầu
toàn bộ lưu lượng định tuyến qua hạ tầng của nhà cung cấp, điều không thể chấp
nhận trong nhiều bối cảnh triển khai đòi hỏi kiểm soát dữ liệu tuyệt đối. Chi phí
cũng là rào cản thực tế: Cloudflare Pro bắt đầu từ 20 USD/tháng mỗi tên miền,
và AWS WAF tính phí theo số lượng WebACL, rule, và triệu lượt yêu cầu, khó
dự đoán và kiểm soát ngân sách với traffic lớn.

Ngoài ra, một số giải pháp như Shadow Daemon hay Naxsi (module Nginx) cũng
tồn tại trong không gian mã nguồn mở nhưng thiếu tính năng quản trị tích hợp,
bộ luật không phong phú và không được duy trì tích cực.

Từ quá trình khảo sát trên, có thể rút ra kết luận: trong không gian giải pháp
WAF hiện tại, tồn tại một khoảng trống đáng kể đối với nhóm nhu cầu cụ thể
sau: \textbf{(i)} triển khai tại chỗ, không phụ thuộc bên thứ ba; \textbf{(ii)} viết
bằng ngôn ngữ hiệu năng cao, có thể xử lý lưu lượng lớn với độ trễ thấp; \textbf{(iii)}
tích hợp một cơ chế phân tích thông minh bổ sung cho bộ luật tĩnh, đặc biệt là để
xử lý các trường hợp mà điểm anomaly nằm trong ``vùng xám'' mà bộ luật không
tự tin đưa ra quyết định; và \textbf{(iv)} cung cấp giao diện quản trị tích hợp, cho
phép đội ngũ bảo mật nội bộ giám sát, phân tích và phản ứng sự cố ngay từ một
điểm duy nhất.

Trên cơ sở đó, đề tài xác định mục tiêu cụ thể như sau: \textit{Xây dựng một
hệ thống tường lửa ứng dụng web (WAF) triển khai tại chỗ, sử dụng ngôn ngữ
lập trình Go, tích hợp hai lớp phát hiện tấn công -- lớp rule engine dựa trên anomaly
scoring và lớp mô hình học máy (DistilBERT fine-tuned) kích hoạt trong vùng xám
-- kèm theo dashboard quản trị real-time, hệ thống xác thực và phân quyền, cơ
chế cảnh báo sự kiện bảo mật và khả năng triển khai đầy đủ qua Docker Compose.}

Phạm vi của đề tài bao gồm:
\begin{itemize}[leftmargin=*]
    \item \textbf{Loại tấn công được bảo vệ}: SQL Injection (SQLi), Cross-Site
    Scripting (XSS), Remote/OS Command Injection (CMDi/RCE), Path Traversal
    và Local File Inclusion (LFI), Server-Side Request Forgery (SSRF), XML
    External Entity (XXE), Log4Shell, Server-Side Template Injection (SSTI),
    NoSQL Injection, Scanner Detection và Rate-based Behavior Analysis. Đây
    là các nhóm tấn công thuộc hoặc liên quan trực tiếp đến OWASP Top 10 (2021)
    và CWE/SANS Top 25.

    \item \textbf{Môi trường triển khai}: máy chủ Linux, container Docker, không
    yêu cầu kết nối Internet sau khi cài đặt. WAF hoạt động theo mô hình
    reverse proxy trong suốt (transparent reverse proxy), không yêu cầu thay
    đổi cấu hình ứng dụng backend.

    \item \textbf{Ngoài phạm vi}: bảo vệ tầng mạng (Layer 3/4), chống DDoS
    volumetric, phân tích lưu lượng mã hóa end-to-end (TLS termination ở
    tầng ngoài), và kiểm tra nội dung response phía backend (chỉ inspect
    request trong phạm vi đề tài này).
\end{itemize}

% ====================================================================
\section{Định hướng giải pháp}
\label{section:1.3}

Từ mục tiêu và phạm vi đã xác định, đề tài đề xuất kiến trúc giải pháp dựa trên
hai nguyên tắc cốt lõi: \textit{hiệu năng không thỏa hiệp ở đường nhanh} và
\textit{độ chính xác tăng cường ở trường hợp mơ hồ}.

\textbf{Lớp thứ nhất -- Rule Engine với Anomaly Scoring.} Toàn bộ logic xử lý
WAF được xây dựng bằng ngôn ngữ Go (Golang). Go là ngôn ngữ biên dịch tĩnh,
có hệ thống goroutine cho phép xử lý đồng thời hàng nghìn kết nối với overhead
tối thiểu, và garbage collector thế hệ mới với thời gian dừng (stop-the-world)
thường dưới một mili giây -- đặc tính quan trọng đối với proxy xử lý lưu lượng
thời gian thực. Thay vì mô hình block-ngay theo từng rule đơn lẻ, đề tài sử dụng
mô hình anomaly scoring: mỗi rule match đóng góp một giá trị điểm vào tổng điểm
nguy cơ của yêu cầu, và quyết định cuối (ALLOW / CHALLENGE / BLOCK) được
đưa ra dựa trên ngưỡng tích lũy. Cách tiếp cận này giúp giảm false positive vì
một yêu cầu bình thường vô tình khớp một rule nhỏ sẽ không bị chặn ngay, nhưng
một chuỗi tín hiệu tấn công dù mỗi cái yếu vẫn có thể tích lũy đủ điểm để kích
hoạt quyết định phòng thủ. Để chống evasion qua encoding, mỗi rule được trang
bị một chuỗi transform (biến đổi tiền xử lý) áp dụng theo thứ tự -- ví dụ
\texttt{URL\_DECODE → LOWERCASE → COMPRESS\_WHITESPACE} -- trước
khi tiến hành so khớp mẫu. Bộ luật của hệ thống gồm 78 rule -- kết hợp tầng
chữ ký tấn công chuẩn và tầng dấu hiệu (\textit{indicator}) tích lũy -- bao phủ
13 nhóm tấn công, được xây dựng theo cấu trúc JSON có phiên bản, dễ mở rộng
và thay thế không cần biên dịch lại.

\textbf{Lớp thứ hai -- Mô hình học máy trong vùng xám.} Quan sát thực tế cho
thấy với nhiều yêu cầu HTTP, bộ luật tĩnh đưa ra điểm anomaly nằm trong khoảng
trung gian -- cao hơn traffic thông thường nhưng chưa đủ để quyết định block dứt
khoát. Để xử lý ``vùng xám'' này, đề tài fine-tune mô hình DistilBERT trên tập
dữ liệu HTTP request có nhãn gồm mười lớp (normal và chín loại tấn công), và
tích hợp mô hình vào hệ thống dưới dạng một microservice Python/FastAPI độc
lập. Khi điểm anomaly của một yêu cầu rơi vào ngưỡng kích hoạt, WAF sẽ gọi
đồng bộ tới ML service với timeout cố định 800 mili giây, nhận phán quyết của
mô hình (nhãn lớp và độ tin cậy), và điều chỉnh điểm anomaly lên hoặc xuống
tùy theo kết quả. Kiến trúc này có hai đặc tính quan trọng: ML không nằm trên
đường nhanh (hot path) của mọi yêu cầu, nên overhead chỉ xuất hiện ở một phần
nhỏ lưu lượng; và hệ thống fail-open (khi ML service không khả dụng, WAF giữ
nguyên điểm rule và tiếp tục xử lý bình thường).

\textbf{Các thành phần bổ trợ.} Bên cạnh hai lớp phát hiện, hệ thống được hoàn
thiện bởi các thành phần sau: \textbf{(i)} Rate Limiting theo thuật toán Token
Bucket per-IP chống tấn công DDoS ứng dụng và brute force; \textbf{(ii)} quản
lý danh sách IP động (blacklist/whitelist); \textbf{(iii)} hệ thống xác thực và
phân quyền dựa trên JWT và bcrypt, với phân quyền theo vai trò Admin/User bảo
vệ các API quản trị; \textbf{(iv)} Dashboard web giám sát và quản trị real-time
cho phép quan sát lưu lượng, xem nhật ký tấn công, quản lý luật và cấu hình hệ
thống từ một giao diện duy nhất; \textbf{(v)} hệ thống cảnh báo sự kiện bảo mật
(Notifier) hỗ trợ gửi thông báo không đồng bộ qua Slack, Email (SMTP) và
Webhook tùy chỉnh; và \textbf{(vi)} toàn bộ hệ thống, gồm WAF, cơ sở dữ liệu
PostgreSQL và ML service, được đóng gói thành một Docker Compose deployment
hoàn chỉnh.

\textbf{Đóng góp chính của đề tài} bao gồm: \textbf{(a)} thiết kế và hiện thực
hóa cơ chế anomaly scoring hai lớp Rule+ML đầu tiên trong không gian mã nguồn
mở Go, với tính chất fail-open và tích hợp không xâm phạm (non-intrusive) vào
hot path; \textbf{(b)} xây dựng bộ luật JSON có cấu trúc với transform chain
chống evasion, kiểm thử và đánh giá trên 42 payload OWASP chuẩn; \textbf{(c)}
xây dựng và đánh giá pipeline fine-tuning DistilBERT cho phân loại HTTP request
10 lớp với cơ chế hard gate đảm bảo chất lượng mô hình trước khi đưa vào triển
khai; và \textbf{(d)} thiết kế interface \texttt{MLPredictor} tách biệt engine và
ML client, cho phép thay thế backend ML mà không thay đổi logic engine.

% ====================================================================
\section{Bố cục đồ án}
\label{section:1.4}

Phần còn lại của báo cáo đồ án tốt nghiệp này được tổ chức như sau.

Chương 2 tập trung khảo sát hiện trạng và phân tích yêu cầu hệ thống. Nội dung
bao gồm khảo sát chi tiết nhu cầu của các nhóm người dùng mục tiêu (Security
Engineer, DevOps, Developer), so sánh mở rộng các giải pháp WAF hiện có, đặc
tả các ca sử dụng quan trọng bằng biểu đồ UML, và xác định đầy đủ yêu cầu chức
năng lẫn phi chức năng làm cơ sở thiết kế cho các chương tiếp theo.

Trong Chương 3, đồ án giới thiệu nền tảng lý thuyết và công nghệ được ứng dụng.
Các chủ đề bao gồm kiến trúc và đặc tính của ngôn ngữ Go liên quan đến xây
dựng proxy hiệu năng cao, mô hình anomaly scoring, cơ sở lý thuyết của kiến trúc
Transformer và quá trình fine-tuning DistilBERT cho bài toán phân loại chuỗi,
thuật toán Token Bucket cho rate limiting, và các công nghệ hạ tầng như PostgreSQL,
JWT, Docker.

Chương 4 trình bày quá trình phân tích thiết kế, triển khai và đánh giá toàn bộ
hệ thống. Đây là chương trọng tâm của đồ án, bao gồm: kiến trúc tổng quan và
luồng xử lý request xuyên suốt các lớp middleware; thiết kế chi tiết rule engine
với cấu trúc dữ liệu, transform chain và cơ chế scoring; tích hợp ML layer và
pipeline huấn luyện mô hình với cơ chế kiểm soát chất lượng; thiết kế cơ sở dữ
liệu; xây dựng dashboard và API quản trị; và trình bày kết quả kiểm thử toàn diện
trên các tập payload tấn công chuẩn, bài kiểm thử tải và kiểm thử bảo mật nội bộ.

Chương 5 tổng hợp và thảo luận sâu về bốn đóng góp nổi bật của đề tài: cơ chế
anomaly scoring hai lớp với ML gray-zone activation; transform chain chống evasion
và tập luật 78 rules; pipeline fine-tuning DistilBERT với hard gate kiểm soát chất
lượng; và thiết kế interface tách biệt engine và ML client phục vụ khả năng kiểm
thử và mở rộng.

Cuối cùng, Chương 6 đúc kết mức độ hoàn thành mục tiêu đề ra, nhìn nhận trực
tiếp các hạn chế còn tồn tại của hệ thống, và đề xuất các hướng phát triển tiếp
theo trong ngắn hạn, trung hạn và dài hạn, bao gồm việc hoàn thiện mô hình ML
phiên bản tiếp theo, mở rộng sang inspection tầng response, và tích hợp cơ chế
học chủ động từ phản hồi của đội ngũ bảo mật.

\bigskip
\noindent\textbf{Kết chương.}
Chương này đã phác thảo toàn cảnh bài toán: sự gia tăng của tấn công tầng ứng
dụng web, hạn chế cố hữu của WAF rule-based và các giải pháp đám mây hiện tại,
khoảng trống còn chưa được lấp đầy trong không gian mã nguồn mở, và định hướng
giải pháp hai lớp mà đề tài theo đuổi. Kiến trúc đề xuất kết hợp rule engine
anomaly scoring viết bằng Go với mô hình DistilBERT fine-tuned kích hoạt có
điều kiện, nhằm đạt được đồng thời hiệu năng cao ở đường nhanh và độ chính xác
tăng cường ở vùng xám. Đây là nền tảng để Chương 2 đi sâu vào khảo sát và
phân tích yêu cầu chi tiết cho từng thành phần của hệ thống.

\end{document}
